\chapter{Use of Generative AI Tools}
\label{app:use_of_generative_ai}
\addcontentsline{toc}{chapter}{Use of Generative AI Tools}

\hspace*{\parindent} The following section details the use of Generative Artificial Intelligence (AI) tools in the course of the development of this project, in accordance with academic integrity policies.


\paragraph{1. Specification of AI-Assisted Components}\hspace*{\parindent}


\hspace*{\parindent} AI assistance was used for the structuring and preparation of the following parts:
\begin{itemize}
    \item Chapter 1: Introduction: flow organization and grammar improvement.
    \item Chapter 2 (Review of Related Literature and Systems): literature search, system identification and comparison, and flow organization.
    \item Chapter 3: Technical Background: flow organization, technical documentation clarity, and grammar improvement.
    \item Chapter 4: Methodology: flow organization, methodological structuring, formatting consistency, and grammar improvement.
    \item In-Text and End-Text Citations: cross-referencing, citation formatting, and validation.
    \item General Manuscript: grammar improvement and vocabulary improvement.
    \item System Design Documentation: flow organization and specification of modules and features.
\end{itemize}


\paragraph{2. Explanation of AI Usage}\hspace*{\parindent}


\hspace*{\parindent} AI tools were used to enhance the quality, clarity, and organization of the manuscript. The AI contribution to Chapter 1: Introduction was flow organization and grammar. The AI contribution to Chapter 2 was literature organization and comparison tables development. The AI contribution to Chapter 3: Technical Background was flow organization, technical documentation clarity, and grammar refinement. The AI contribution to Chapter 4: Methodology was flow organization, methodological structuring, formatting consistency, and clarity in presenting system implementation and evaluation details. AI tools helped with consistent cross-referencing, citation formatting, and grammar. Additionally, AI was used to generate design concepts and refine the KABISIG logo. Usage of AI tools was restricted only to editorial, organizational, and design purposes. All research conclusions, analyses, interpretations, designs of the system, and writing of the manuscript are the intellectual property of the researchers.\vspace{5mm}


\paragraph{3. Justification of Necessity}\hspace*{\parindent}


\hspace*{\parindent} AI tools were required to increase clarity, cohesion, and quality of the manuscript. They helped with organizing a large amount of information in Chapter 2 and consistent citation formatting, as well as structuring technical explanations and workflows in Chapters 3: Technical Background and Chapter 4: Methodology, thus reducing errors and making it possible for researchers to focus more on analysis and design. AI was also used to generate logo concepts, streamlining the design process and providing visual direction for the project's branding.\vspace{5mm}


\paragraph{4. Affirmation of Understanding}\hspace*{\parindent}


\hspace*{\parindent} The authors recognize, understand, and are responsible for all AI-assisted parts in this manuscript. AI tools were used to assist with literature search, organizing of flow, editing, formatting, and logo design. All sources of literature that were found using searches were checked manually. The authors are responsible for all aspects of this work, including originality and integrity of this paper.